\documentclass[10pt]{article}
\usepackage[T1]{fontenc}
\usepackage[utf8]{inputenc}
\usepackage[french]{babel}
\usepackage{amsmath,amssymb,mathrsfs,mathpazo}
\usepackage{geometry}
\geometry{a4paper,margin=2cm}
\usepackage{graphicx,array,tabularx,tabularray,colortbl,multirow,multicol,booktabs,enumitem,pifont,ulem,fancybox,diagbox,cancel}
\usepackage{pstricks,pstricks-add,pst-plot,pst-tree,pst-node,pst-text,pst-func,pst-math,pst-bspline,pst-3dplot,pst-eucl,pst-solides3d}
\let\pscalculate\undefined
\usepackage{pst-bezier}
\usepackage{tikz,pgfplots,tkz-tab}
\usetikzlibrary{arrows,arrows.meta,patterns,calc,positioning,trees,decorations.markings,shapes.geometric,shapes.misc,shapes.arrows,intersections,decorations.pathmorphing,decorations.pathreplacing,backgrounds,shadows,babel}
\pgfplotsset{compat=1.18}
\usepgfplotslibrary{fillbetween}
\usepackage[np]{numprint}
\usepackage{listings,makecell,pas-tableur,twoopt}
\let\barre\undefined
\usepackage[pstricks]{tablvar}
\definecolor{darkgreen}{RGB}{0,100,0}
\providecommand{\red}{\color{red}}
\providecommand{\blue}{\color{blue}}
\providecommand{\green}{\color{green}}
\providecommand{\black}{\color{black}}
\providecommand{\white}{\color{white}}
\providecommand{\gray}{\color{gray}}
\providecommand{\cyan}{\color{cyan}}
\providecommand{\magenta}{\color{magenta}}
\providecommand{\yellow}{\color{yellow}}
\providecommand{\violet}{\color{violet}}
\providecommand{\orange}{\color{orange}}
\providecommand{\euro}{\texteuro}
\providecommand{\EUR}{\texteuro}
\providecommand{\ieme}{\textsuperscript{e}}
\providecommand{\degres}{\ensuremath{^\circ}}
\providecommand{\degree}{\ensuremath{^\circ}}
\providecommand{\vv}[1]{\overrightarrow{#1}}
\providecommand{\Coord}[1]{#1}
\providecommand{\decosix}{}
\providecommand{\starredbullet}{\ensuremath{\bullet}}
\pagestyle{empty}
\setlength{\parindent}{0pt}
\usepackage{adjustbox}\DeclareUnicodeCharacter{2236}{:}\DeclareUnicodeCharacter{2190}{\ensuremath{\leftarrow}}
\providecommand{\R}{\mathbb{R}}
\providecommand{\N}{\mathbb{N}}
\providecommand{\Z}{\mathbb{Z}}
\providecommand{\Q}{\mathbb{Q}}
\providecommand{\D}{\mathbb{D}}
\providecommand{\C}{\mathbb{C}}
\providecommand{\e}{{\mathrm{e}}}
\providecommand{\vect}[1]{\overrightarrow{#1}}
\providecommand{\vectt}[1]{\overrightarrow{\mathrm{#1}}}
\providecommand{\vvt}[1]{\overrightarrow{\mathrm{#1}}}
\providecommand{\np}[2][]{{#2~#1}}
\providecommand{\vv}[1]{{\overrightarrow{#1}}}
\renewcommand{\d}{{\mathrm{d}}}
\providecommand{\degres}{{^\circ}}
\providecommand{\degre}{{^\circ}}
\providecommand{\blue}{\color{blue}}
\providecommand{\red}{\color{red}}
\providecommand{\enskip}{\quad}
\providecommand{\llbracket}{\lbrack\!\lbrack}
\providecommand{\rrbracket}{\rbrack\!\rbrack}
\providecommand{\euro}{{\text{€}}}
\providecommand{\no}{{\text{n°}}}
\providecommand{\interval}[2]{{[#1;#2]}}
\providecommand{\intervFO}[2]{{[#1;#2[}}
\providecommand{\dotfill}{\ldots}

\providecommand{\ds}{\displaystyle}
\providecommand{\textding}[1]{#1}
\providecommand{\Oij}{\(\mathrm{O};\vec{\imath},\vec{\jmath}\)}
\providecommand{\Oijk}{\(\mathrm{O};\vec{\imath},\vec{\jmath},\vec{k}\)}
\providecommand{\Ouv}{\(\mathrm{O};\vec{u},\vec{v}\)}
\providecommand{\intervFF}[2]{\ensuremath{[#1;#2]}}
\renewcommand{\euro}{\ifmmode\text{\texteuro}\else\texteuro\fi}
\DeclareUnicodeCharacter{00B0}{\ifmmode^{\circ}\else\textdegree\fi}
\DeclareUnicodeCharacter{20AC}{\ifmmode\text{\texteuro}\else\texteuro\fi}
% Énoncés : Éducation nationale. Transcription et corrigés : APMEP, Denis Vergès et contributeurs. Découpage et classement : CoopMaths, Sébastien Lozano, Éric Elter et contributeurs. Adaptation de la présentation : Maths Explorer.
% CC BY-SA : https://coopmaths.fr/www/about/
% https://coopmaths.fr/alea/static/bac/2025/tex/bac_2025_05_sujet2_ameriquenord_2.tex
\begin{document}
\noindent\textit{APMEP / CoopMaths. CC BY-SA. Présentation adaptée par Maths Explorer.}\par\medskip
\medskip

Un des objectifs de cet exercice est de déterminer une approximation du nombre réel $\ln(2)$, en utilisant une des méthodes du mathématicien anglais Henry Briggs au XVI\up{e} siècle.

\bigskip

On désigne par $\left(u_{n}\right)$ la suite définie par :

\[u_{0}=2 \quad \text { et, pour tout entier naturel } n,\quad  u_{n+1}=\sqrt{u_{n}}\]

\medskip

\textbf{Partie A}

\medskip

\begin{enumerate}
\item 
	\begin{enumerate}
		\item Donner la valeur exacte de $u_{1}$ et de $u_{2}$.
		\item Émettre une conjecture, à l'aide de la calculatrice, sur le sens de variation et la limite éventuelle de la suite.
	\end{enumerate}
\item 
	\begin{enumerate}
		\item Montrer par récurrence que pour tout entier naturel $n$,\quad $1 \leqslant u_{n+1} \leqslant u_{n}$.
		\item En déduire que la suite $\left(u_{n}\right)$ est convergente.
		\item Résoudre dans l'intervalle $[0~;~+\infty[$ l'équation $\sqrt{x} =x$.
		\item Déterminer, en justifiant, la limite de la suite $\left(u_{n}\right)$.
	\end{enumerate}
\end{enumerate}

\medskip

\textbf{Partie B}

\medskip
	On désigne par $\left(v_{n}\right)$ la suite définie pour tout entier naturel $n$ par $v_{n}=\ln \left(u_{n}\right)$.

\medskip

\begin{enumerate}
\item 
	\begin{enumerate}
		\item Démontrer que la suite $\left(v_{n}\right)$ est géométrique de raison $\dfrac{1}{2}$.
		\item Exprimer $v_{n}$ en fonction de $n$, pour tout entier naturel $n$.
		\item En déduire que, pour tout entier naturel $n$,\quad $\ln (2)=2^{n} \ln \left(u_{n}\right)$.
	\end{enumerate}
\item On a tracé ci-dessous dans un repère orthonormé la courbe $\mathscr{C}$ de la fonction ln et la tangente T à la courbe $\mathscr{C}$ au point d'abscisse 1.

Une équation de la droite T est $y = x - 1$.

Les points $\mathrm{A}_{0}$, $\mathrm{A}_{1}$, $\mathrm{A}_{2}$ ont pour abscisses respectives $u_{0}$, $u_{1}$ et $u_{2}$ et pour ordonnée~0.

\begin{center}













\psset{xunit=4.5cm, yunit=4.5cm,arrowsize=2pt 3}
\begin{pspicture*}(-0.2,-0.2)(2.2,1.2)
\psaxes[linewidth=0.8pt,labels=none,Dx=1,Dy=0.5]{->}(0,0)(-0.2,-0.2)(2.2,1.2)

\psline[linewidth=0.5pt](1,-0.05)(1,0.05)
\rput[t](1,-0.1){\footnotesize 1}
\psline[linewidth=0.5pt](2,-0.05)(2,0.05)
\rput[t](2,-0.1){\footnotesize 2}
\psline[linewidth=0.5pt](-0.05,0.5)(0.05,0.5)
\rput[r](-0.1,0.5){\footnotesize 0{,}5}
\psline[linewidth=0.5pt](-0.05,1)(0.05,1)
\rput[r](-0.1,1){\footnotesize 1}
\rput[tr](-0.05,-0.05){\footnotesize 0}

\psplot[linewidth=1.25pt,linecolor=red,plotpoints=2000]{0.6}{2.2}{x ln}
\rput[l](1.9,0.58){\red $\mathscr{C}$}
\psplot[linewidth=1pt, linecolor=blue, plotpoints=10]{0.6}{2.2}{x 1 sub}
\rput[l](1.9,0.8){\blue T}

\psdot[dotsize=2pt, linecolor=black](2,0)
\rput[b](2,0.05){A$_0$}
\psdot[dotsize=2pt, linecolor=black](1.414,0)
\rput[b](1.414,0.05){A$_1$}
\psdot[dotsize=2pt, linecolor=black](1.189,0)
\rput[b](1.189,0.05){A$_2$}
\end{pspicture*}
\end{center}

On décide de prendre $x - 1$ comme approximation de $\ln (x)$ lorsque $x$ appartient à l'intervalle $]0,99~;~1,01[$.

	\begin{enumerate}
		\item Déterminer à l'aide de la calculatrice le plus petit entier naturel $k$ tel que $u_{k}$ appartienne à l'intervalle $]0,99 ~;~1,01[$ et donner une valeur approchée de $u_{k}$ à $10^{-5}$ près.
		\item En déduire une approximation de $\ln \left(u_{k}\right)$.
\item Déduire des questions \textbf{1. c.} et \textbf{2. b.} de la \textbf{partie B} une approximation de $\ln (2)$.
	\end{enumerate}

\item On généralise la méthode précédente à tout réel $a$ strictement supérieur à 1.

Recopier et compléter l'algorithme ci-dessous afin que l'appel \texttt{Briggs(a)} renvoie une approximation de $\ln (a)$.

On rappelle que l'instruction en langage Python \texttt{sqrt(a)} correspond à $\sqrt{a}$.

\begin{center}
\begin{ttfamily}
\begin{tabularx}{8cm}{|X|}\hline
\textbf{from} math \textbf{import}*\\
\textbf{def} Briggs(a):\\
\quad n = 0\\
\quad \textbf{while} a >= 1.01:\\
	\quad \quad a = \textbf{sqrt}(a)\\
	\quad \quad n = n+1\\
	\quad L =\dots\\
	\textbf{return} L\\ \hline
	\end{tabularx}
	\end{ttfamily}
\end{center}
\end{enumerate}



\bigskip

\end{document}
